\documentclass[11pt,a4paper]{article}
\usepackage[margin=27mm]{geometry}
\usepackage{fontspec,unicode-math,amsmath,amsthm,booktabs,longtable,hyperref}
\setmainfont{STIXGeneral}\setmathfont{STIX Two Math}\hypersetup{hidelinks}
\newtheorem{theorem}{Theorem}\newtheorem{lemma}[theorem]{Lemma}\newtheorem{proposition}[theorem]{Proposition}
\newcommand{\Cl}{\operatorname{Cl}}\newcommand{\Aut}{\operatorname{Aut}}\newcommand{\Cen}{\operatorname{C}}
\title{Conjugacy tests for quotients of a locally projective polytope}
\author{A reproof of the Hartley--Leemans erratum}\date{}
\begin{document}\maketitle
The two-page erratum of Hartley and Leemans \cite{erratum} removes eight entries from the quotient census in \cite{original}. We recover every corrected census number, identify the eight deleted rows from the printed data, and match all surviving rows and their subgroup relations. The general observation is that a condition on every conjugate of a subgroup can be checked on element conjugacy classes. Thus no conjugate is omitted through sampling. The exhaustive finite calculation is a required part of this proof.

\section{A complete conjugacy test}
Let $G$ be finite, let $B\subseteq G$ be a forbidden set, and let $N\leq G$. Our convention is $x^g=g^{-1}xg$.
\begin{lemma}\label{avoid}
The following conditions are equivalent:
\[
 N^g\cap B=\varnothing\quad(g\in G),\qquad
 N\cap\bigcup_{b\in B}\Cl_G(b)=\varnothing.
\]
For the second condition it suffices to test one representative of every conjugacy class of elements of $N$. The admissible subgroups are closed under conjugation and passage to subgroups.
\end{lemma}
\begin{proof}
An element $b$ lies in $N^g\cap B$ exactly when $b=x^g$ for some $x\in N$ and $b\in B$. This is precisely membership of $x$ in a $G$-class meeting $B$. Conjugation within $N$ does not change its $G$-class, proving the representative test. Both closure assertions follow directly from the first condition.
\end{proof}

Suppose that $G$ is the marked group of a regular polytope, that $K$ is the stabilizer of a facet, and that $\omega$ is a central involution of $K$. Assume the quotient under consideration has flags indexed by the cosets of $N$, and that every intersection $N^g\cap K$ is either $1$ or $\langle\omega\rangle$. These assumptions concern the actual quotient and local sections, not just an arbitrary permutation action.
\begin{proposition}\label{facets}
Put $m=|N\cap\Cl_G(\omega)|$ and $c=|\Cen_G(\omega)|$. If $h$ and $d$ count the facet orbits with stabilizer of order two and one, respectively, then
\[
 h=\frac{2cm}{|K||N|},\qquad
 d=\frac{|G|}{|K||N|}-\frac h2.
\]
Moreover the full automorphism group of the quotient is $N_G(N)/N$.
\end{proposition}
\begin{proof}
Facets are cosets of $K$. A coset $Kg$ has an order-two stabilizer in $N$ exactly when $\omega^g\in N$. There are $cm$ choices of $g$ satisfying this condition: each of the $m$ possible conjugates has $c$ preimages. Centrality of $\omega$ in $K$ makes this condition independent of the representative of $Kg$, so there are $cm/|K|$ such facets. Their $N$-orbits have size $|N|/2$, which gives the formula for $h$. Counting all $|G|/|K|$ facets by their orbit sizes gives the formula for $d$.

On flags the marked generators act by adjacency. A type-preserving automorphism of the quotient commutes with these generators and hence with their generated transitive action. An equivariant bijection of $N\backslash G$ is determined by the image of $N$. That image is $Ng$ with the same stabilizer precisely when $g$ normalizes $N$. The resulting permutations are $N_G(N)/N$. Conversely, these permutations preserve every flag adjacency and therefore the quotient faces and incidence. This proves the automorphism assertion.
\end{proof}

\begin{lemma}\label{maxima}
Let a conjugacy-invariant collection of subgroups of a finite group be closed under subgroups. Form a directed graph on its subgroup conjugacy classes by recording, for each representative, the ambient conjugacy classes of all its maximal proper subgroups. A vertex is maximal in the collection exactly when it has no incoming edge.
\end{lemma}
\begin{proof}
An incoming edge exhibits a proper containing member. Conversely, if $N$ is properly contained in a member, choose a containing member minimal above $N$. Finiteness makes $N$ a maximal subgroup of that member. Conjugating the containing member to the chosen representative gives the required recorded edge. Repeated edges are counted only once.
\end{proof}

\section{The exact polytope and the retained geometric inputs}
Write $P$ for the universal regular polytope with dodecahedral facets and hemi-icosahedral vertex figures. We retain the following results of the original article, not its erroneous census:
\begin{enumerate}
\item Its Theorems 3.1--3.2 identify its marked group as
\[
 W=[5,3,5]/\langle\!\langle(s_1s_2s_3)^5\rangle\!\rangle
   \cong J_1\times\operatorname{PSL}_2(19),\quad |W|=600415200.
\]
The marking consists of involutions $s_0,s_1,s_2,s_3$; nonadjacent generators commute, and adjacent product orders are $5,3,5$.
\item With $H_0=\langle s_1,s_2,s_3\rangle$, $H_3=\langle s_0,s_1,s_2\rangle$ and $\omega=(s_0s_1s_2)^5$, the quotient criterion stated in its Section 4 is
\begin{equation}\label{criterion}
 N\text{ is semisparse}\quad\Longleftrightarrow\quad
 N^w\cap H_0H_3\subseteq\{1,\omega\}\quad(w\in W).
\end{equation}
This is the source's application of Hartley's endpoint criterion: the hemi-icosahedron has no proper quotient, and the permitted dodecahedral facet intersections are $1$ and its antipodal central subgroup. Its quotient correspondence identifies isomorphism classes of the quotient polytopes with these subgroup conjugacy classes.
\end{enumerate}
These are substantial retained marked-group and geometric inputs. We do not reprove universal finiteness, the general quotient correspondence, or the general semisparse theorem here, and do not claim recovery of the separate original L095. In particular, no assertion that every subgroup with a polytopal orbit poset is semisparse is needed. The quotient convention is the one in \eqref{criterion} and the cited correspondence.

The exact permutation model is reconstructed as follows. The existing accepted Janko marking search from the common chamber-operations calculation is reused, with the generator order reversed. Its only library input is the named TomLib permutation model of $J_1$. The other factor uses the four literal permutations on twenty points displayed in Section 3 of \cite{original}. The direct-product images of corresponding markers generate the whole product; the program checks every defining relator, its order and the two endpoint orders $60,120$. Combined with the retained presentation theorem, the surjection from the presented $W$ onto this model is an isomorphism. The model is therefore marked, not merely an abstract group of the right order. The complete reused search and the twenty-point permutations are included in \texttt{ambient\_marking.g}; no old manuscript is a hidden proof premise.

Apply Lemma~\ref{avoid} to
\[
 B=H_0H_3\setminus\{1,\omega\}.
\]
The product has $1200$ elements, meeting $47$ of the $180$ element conjugacy classes of $W$ after the two permitted elements are removed. Every representative of every element class in each subgroup is tested against this full forbidden-class set. For a surviving subgroup, \eqref{criterion} also guarantees exactly the local alternatives required by Proposition~\ref{facets}. All vertex figures remain hemi-icosahedra; its $d$ facets are dodecahedra and its $h$ facets hemi-dodecahedra.

\section{The corrected census and the historical labels}
\begin{theorem}\label{census}
The group $W$ has $137$ conjugacy classes of semisparse subgroups, of which $26$ are maximal. The corresponding quotient polytopes comprise $65$ with the original facets and vertex figures, one with only hemi-dodecahedral facets, and $71$ with mixed facets. Thus $66$ are section regular. The eight rows removed in \cite{erratum} are exactly
\[
 13,\ 70,\ 78,\ 90,\ 91,\ 95,\ 111,\ 127.
\]
\end{theorem}
\begin{proof}
The required program uses GAP's algorithm enumerating every subgroup conjugacy class, yielding $1262$ representatives. It does not read the authors' old subgroup file, a semisparse list, or a list of quotient polytopes. Lemma~\ref{avoid} filters this exhaustive list to $137$ classes. Proposition~\ref{facets} gives $65$ rows with $h=0$, one with $d=0$, and $71$ with $d,h>0$. The two homogeneous facet cases are precisely section regular here, since the vertex figure is fixed and all lower sections are those of the two stated regular facets. The program separately enumerates maximal subgroups of every surviving representative. The resulting graph on ambient subgroup classes has $349$ distinct edges and, by Lemma~\ref{maxima}, $26$ maximal vertices.

For the source labels, all thirty generating sets in the original Table 1 are evaluated as literal words in this same marking. Exactly the four labelled $13,70,95,127$ fail the test; the other twenty-six are exactly the computed maximal classes. The four remaining deletions are identified from the printed inclusion rows as follows:
\begin{center}\begin{tabular}{cll}\toprule
Row&Printed subgroup&Identification in the literal parent\\\midrule
78&$C_{45}$&the unique subgroup class of order $45$ in row 13\\
111&$C_{15}$&the unique subgroup class of order $15$ in row 13\\
90&$C_{30}$&the unique cyclic subgroup class of order $30$ in row 70\\
91&$C_3\times D_{10}$&the order-$30$ nonabelian child of row 70\\
&&with normalizer quotient of order $12$\\\bottomrule\end{tabular}\end{center}
Here $D_{10}$ has order ten. The other nonabelian order-$30$ child has normalizer quotient of order $36$ and is not silently identified with row 91. For each of the eight rows the output gives actual permutations $x,q,b$ with
\[
 x\in N,\quad q\in W,\quad x^q=b\in H_0H_3\setminus\{1,\omega\}.
\]
These are explicit failures of the required universal conjugate condition.

Finally, the $137$ surviving historical rows, using all their printed maximal-subgroup links and facet counts, are matched bijectively to the computed $137$-vertex, $349$-edge graph. The thirty literal generating sets anchor the maximal rows; every surviving row is below a surviving maximum. The exact graph isomorphism, orders and both facet counts are checked, rather than only the grand total. All $137$ normalizer quotients also agree with the printed automorphism columns, after elementary direct-product identities. These include $D_{12}\cong C_2\times S_3$, $D_{20}\cong C_2\times D_{10}$, splitting coprime cyclic factors, and $\operatorname{GL}_2(4)\cong C_3\times A_5$: in the last identity the scalar center has order three, meets $\operatorname{SL}_2(4)\cong A_5$ trivially, and together these subgroups have the full order $180$. The complete original row set minus these survivors is exactly the displayed eight labels. The required finite output and the group/quotient premises above therefore establish all assertions.
\end{proof}

\paragraph{Edition crosswalk.}
Our complete original source is the fifteen-page 2003 arXiv edition of \cite{original}. Its Table 1 labels an order-$70$ subgroup $C_5\times D_{14}$ as ``58'', whereas Tables 2 and 5 label the order-$76$ group $C_2^2\times C_{19}$ as 58 and this order-$70$ group as 63. The literal word is kept with its Table 1 label; the crosswalk to the full subgroup tables is recorded explicitly as 63. Two further column discrepancies occur on its page 14: row 21 prints $2^2\times(19:3)$ although its facet count forces order $456$, and row 27 prints $2\times19:3$ although its facet count forces order $342$. The corresponding exact subgroup computations give $2^3\times(19:3)$ and $2\times(19:9)$ respectively. All literal extracted entries remain in the data; the comparison does not replace them silently. These are inconsistencies of the accessible edition. We have not inspected the twelve-page journal edition and do not attribute these column discrepancies to it. They do not change the eight labels in the 2006 erratum.

\paragraph{Evidence and attribution.}
The correction targets and all historical words and tables are Hartley--Leemans data. The present contribution is the common conjugacy-avoidance reduction, the explicit orbit count and maximality argument, and a fresh exhaustive implementation with a complete historical-row match. GAP 4.15.1's exact finite-group algorithms and the installed TomLib named $J_1$ model are retained computational premises. Their versions and file hashes are bound by the runner. Every subgroup representative, each rejection class, the complete surviving containment graph, facet counts, normalizer quotients and eight literal failure witnesses are saved. The programs are required evidence, not sampled validation. Negative controls change the forbidden classes, omit an ambient factor, alter a source label, remove a local facet type, and corrupt the subgroup graph. No Lean or end-to-end formal-verification claim is made. Timing history and the authors' diagnosis of their former implementation are historical statements, not newly proved assertions. No priority or new-conjecture claim is made.

\begin{thebibliography}{9}
\bibitem{erratum} M. I. Hartley and D. Leemans, \emph{Quotients of a universal locally projective polytope}, erratum, Math. Z. \textbf{253} (2006), 433--434, doi:10.1007/s00209-005-0872-z.
\bibitem{original} M. I. Hartley and D. Leemans, \emph{Quotients of a universal locally projective polytope of type $\{5,3,5\}$}, arXiv:math/0307075v1 (2003), fifteen pages; journal reference Math. Z. \textbf{247} (2004), 663--674. The edition distinction is preserved above.
\bibitem{semisparse} M. I. Hartley, \emph{Simpler Conditions for Semisparse Subgroups}, cited as submitted in reference [7] of the 2003 edition of \cite{original}. The precise retained endpoint criterion is read through its explicit statement and application in \cite{original}; no fresh complete reading of the later journal paper is claimed.
\bibitem{reuse} Existing common theory for chamber operations, complete Janko marking program \texttt{phase40\_compress\_operations\_check.g}. Only its marking search is reused and printed in the current required program; the earlier geometric census is not an input.
\end{thebibliography}
\end{document}
